\section{Modernes Malware Design: Umgegung des Microsoft Defender}
\subsection*{Dokumentieren Sie Ihr Vorgehen, um den MS Defender mit einer Malware zu umgehen. Dies simuliert eine Social-Engineering-Attacke, bei der der Nutzer die Malware bspw. via Office-Dokument nachlädt (die Malware wird also von einem externen Server implizit nachgeladen, wenn das Dokument zum Editieren geöffnet wird, bspw. via Word-Makro). Sie müssen hierbei die Malware nicht selbst entwickeln, es geht hier nur um das theoretische Vorgehen und die einzelnen Schritte.}
Als Grundlage dient das Paper „Bypass Antivirus Dynamic Analysis“ von Emeric Nasi.\\
Grundlegend basieren Umgehungen von Antivirenprogrammen auf zwei sicherzustellenden Schritten. Der Code, welcher möglicherweise als gefährlich eingestuft werden könnte, muss versteckt werden, meist durch Verschlüsselung.\\
Ebenso muss der Decryption Stub so erstellt werden, dass er nicht als Virus erkannt wird und die Emulations- bzw. Sandboxing-Phase umgeht.\
Um zu verstehen, wie genau das funktioniert, müssen zunächst die Grundlagen betrachtet werden, anhand derer ein Antivirusprogramm erkennt, ob es sich um Malware handelt. Zum einen gibt es die Static Signature Analysis. Wenn eine Malware erkannt wird, bekommt sie eine Signatur und wird in einer großen Datenbank gespeichert, mit der zukünftige Programme beim Analysieren abgeglichen werden. Dadurch kann bekannte Malware verhindert werden, sofern sie nicht großartig modifiziert wurde. Dennoch bringt diese Technik nichts gegen neue, bislang unbekannte Malware. Die zweite Technik eines Antivirusprogramm ist die Static Heuristic Analysis. Hierbei wird der Code nach bekannten Mustern durchsucht, die auf Malware hindeuten könnten. Damit kann auch noch nicht bekannte Malware erkannt werden. Allerdings entstehen dabei auch False Positives.\
Die letzte Methode ist die Dynamic Analysis. Hierbei wird der zu scannende Code in einer virtuellen Umgebung ausgeführt. Diese Methode ist populär, da viel Software verschlüsselt ist. Das Ziel des Antivirusprogramms besteht darin, dass der Code durch die Ausführung entschlüsselt wird und somit mit den anderen Verfahren der „neue“ bzw. eigentliche Code analysiert werden kann und nicht nur die Verschleierung. Wie kann man dies nun aushebeln?\\
Wie bereits erwähnt, ist das Verschlüsseln eines der gängigen Mittel, um dem Antivirusprogramm Probleme zu bereiten. Zusätzlich benötigt man einen nicht erkannten Self-Decryption-Mechanismus, und im besten Fall kann man das Antivirusprogramm daran hindern, den Decryption Stub auszuführen. Eine weitere wichtige Limitation, die man kennen muss, um Antivirenprogramme zu täuschen, ist, dass Scans sehr schnell durchgeführt werden müssen und somit die Anzahl der Operationen für jeden Scan limitiert ist. Gleichzeitig ist die simulierte Umgebung nicht identisch mit der realen Maschinenumgebung, was ebenfalls von Malware erkannt werden kann.\
Es gibt also verschiedene Ansätze, Malware so zu verschleiern, dass sie nicht vom Antivirusprogramm erkannt wird. Zum theoretischen Vorgehen in unserem Szenario:\\
Wir haben die Ausgangssituation, dass unser Target auf unseren Anhang klickt und versucht, das Dokument zu bearbeiten. Damit wird im Hintergrund die Malware heruntergeladen und vom MS Defender geprüft. Grundlegend brauchen wir nun, wie bereits vorher auch schon erläutert:
\begin{lstlisting}[caption={Grundbaustein}]
	int main()
	decryptCodeSection();
	startShellCode();
\end{lstlisting}
Dies alleine reicht aber noch nicht, da der Code beim Ausführen entschlüsselt wird und somit erkannt wird, dass eine Meterpreter-Shell geöffnet werden soll.\\
Daher bedienen wir uns an dem Beispiel aus dem Paper, bei dem wir vor dem Entschlüsseln und Starten des Shellcodes sehr viel Speicher reservieren und wieder freigeben. Dabei nutzen wir aus, dass der Defender den Scan abbricht, wenn zu viele Ressourcen beansprucht werden. So schaffen wir es zu umgehen, dass unser Code geflaggt wird. Eine sehr simple, aber effektive Methode.\\
Damit wäre das Ziel wahrscheinlich erreicht. Trotzdem wird folgend noch eine weitere Methode benannt, wie es klappen könnte, die ebenfalls in dem Paper vorgestellt wurde.\\
Die „Knowing your enemy“-Methode, mit der wir prüfen können, ob wir uns noch in der Sandbox-Umgebung des MS Defenders befinden. Dies erfordert jedoch etwas erweiterte Informationen über das System des Gegenübers. Da es sich hier in der Aufgabe scheinbar um Social Engineering handelt, könnten diese Informationen bereits vorliegen. Der Trick besteht darin, zu versuchen, spezifische Dateien des Benutzeraccounts auszulesen, und nur dann, wenn dies erfolgreich ist, den Shellcode zu entschlüsseln und auszuführen. In der Sandbox-Umgebung existiert dieser spezifische Nutzer bzw. die Datei in der Regel nicht. Somit wird der Shellcode in der Sandbox-Umgebung nicht entschlüsselt und kann folglich auch nicht geflaggt werden.